\documentclass[11pt]{article}
\usepackage[T1]{fontenc}
\usepackage{textcomp}
\usepackage[margin=0.75in]{geometry}
\usepackage{amsmath,amssymb,amsthm}
\usepackage{array,booktabs}
\usepackage{hyperref}
\setlength{\parindent}{0pt}
\newtheorem{theorem}{Theorem}
\theoremstyle{definition}
\newtheorem{definition}{Definition}
\title{Flow Proof Artifact: List-rev-double-cons}
\author{Generated by \texttt{flow doc proof}}
\date{Flow Proof Book}
\begin{document}
\maketitle
Reversing a two-element cons list swaps endpoints.\\[0.5em]
\textbf{Source.} church — https://en.wikipedia.org/wiki/List\_(abstract\_data\_type)\\[0.5em]
\bigskip
\noindent{\large \textbf{Derived fact 1.} \textit{rev(cons(x, cons(y, nil))) = cons(y, cons(x, nil))} \hfill List \cdot reverse \cdot reverse of pair cons swaps endpoints}\par
\smallskip\noindent\textit{Coordinate: List \cdot reverse \cdot reverse of pair cons swaps endpoints}\par
\smallskip\noindent\textit{Source: church}\par
\smallskip\noindent\textit{Built on: reverse of cons appends reversed tail, for reverse on List, singleton reverse is identity, for reverse on List}\par
\medskip
\noindent\textbf{Goal.} rev(cons(x, cons(y, nil))) = cons(y, cons(x, nil))\par
\noindent$\displaystyle \forall x \in \mathbb{N} \forall y \in \mathbb{N}\quad rev(cons(x, cons(y, nil))) = cons(y, cons(x, nil))$\par
\medskip
\noindent
\begin{tabular}{@{} >{\raggedright\arraybackslash}p{0.06\textwidth} >{\raggedright\arraybackslash}p{0.47\textwidth} >{\raggedleft\arraybackslash}p{0.41\textwidth} @{}}
\textbf{\#} & \textbf{Proof} & \textbf{Mathematics} \\
\midrule
\textbf{1.} & We prove that reverse of pair cons swaps endpoints for reverse on List. & \\[0.45em]
\textbf{2.} & We invoke the derived fact governing reverse on List: reverse of cons appends reversed tail, for reverse on List (instantiated for x, cons(y, nil)). & \\[0.45em]
\textbf{3.} & We invoke the derived fact governing reverse on List: singleton reverse is identity, for reverse on List (instantiated for y). & \\[0.45em]
\textbf{4.} & From step 2 and step 3, this implies rev(cons(x, cons(y, nil))) equals cons(y, cons(x, nil)). Hence proven. & $\displaystyle rev(cons(x, cons(y, nil))) = cons(y, cons(x, nil))$ \\[0.45em]
\bottomrule
\end{tabular}
\label{thm:List-reverse-reverse_of_pair_cons_swaps_endpoints}
\par\medskip\hrule
\end{document}
